\begin{enumerate}
\item T, A et M sont alignés dans cet ordre donc

$\mathrm{TA} + \mathrm{AM} = \mathrm{TM}$ d'où $\mathrm{AM} = \mathrm{TM} - \mathrm{TA}$; or $7,7 - 2,2 = 5,5$ donc $\mathrm{AM} =  \np[m]{5,5}$.

\item Le triangle BAM est rectangle en A ; on lui applique le théorème de Pythagore.

$\mathrm{MB}^2 = \mathrm{AM}^2+\mathrm{AB}^2$ \quad donc\quad  $7,3^2  = \mathrm{AB}^2 + \mathrm{AM}^2 $ d'où AB$^2 = 7,3^2 - 5,5^2 = 53,29 - 20,25  = 23,04$.

$\mathrm{AB}>0$ \quad et \quad $\sqrt{23,04}=4,8$ \quad donc \quad  $\mathrm{AB} = \np[m]{4,8}$.

\emph{Remarque} : on pouvait voir que $a^2 - b^2 = (a + b)(a - b)$ et donc que
$7,3^2 - 5,5^2 = (7,3 + 5,5)(7,3 - 5,5) = 12,8 \times 1,8 = 6,4 \times 2 \times 3,6 = 64 \times \times 36 \times 0,01 = 8^2 \times 6^2 \times 0,1^2= (8 \times 6 \times 0,1)^2 = 4,8^2$ : on retrouve $\mathrm{AB} = \np[m]{4,8}$.

\item On sait que  dans le triangle BAM  rectangle en A on a par définition :

$\cos\left(\widehat{\text{ABM}}\right) = \dfrac{\text{longueur du côté adjacent}}{\text{longueur de l'hypoténuse}} = \dfrac{\text{AB}}{\text{MB}} = \dfrac{4,8}{7,3}  \approx \np{0,6575}$.

La calculatrice donne : $\widehat{\text{ABM}}\approx 48,9\degres$ soit environ 49\degres{} au degré près.

	\item \begin{itemize}
		\item E, A et B sont alignés dans cet ordre ;
		\item T, A et M sont alignés dans cet ordre ;
		\item $\dfrac{\mathrm{AE}}{\mathrm{AB}} = \dfrac{1,92}{4,8} = \dfrac{192}{480} = \dfrac{48}{120} = \dfrac{4}{10} = \dfrac25$ ; et $\dfrac{\mathrm{AT}}{\mathrm{AM}} = \dfrac{2,2}{5,5} = \dfrac{22}{55} = \dfrac25$.
	\end{itemize}

On a donc $\dfrac{\mathrm{AE}}{\mathrm{AB}} = \dfrac{\mathrm{AT}}{\mathrm{AM}}$ et d'après
 la réciproque du théorème de Thalès, les droites (TE) et (BM) sont parallèles.
	\item Si le triangle ABM est l'image du triangle AET par une homothétie de centre A, alors le côté homologue de [AT] (dans le triangle de départ) est [AM] (dans le triangle après l'homothétie). Le rapport de l'homothétie est donc $2,5$ ou $- 2,5$ car $\dfrac{\text{AM}}{\text{AT}} = \dfrac{5,5}{2,2} = \dfrac{55}{22} = \dfrac52 = 2,5$ et comme A (centre de l'homothétie est situé entre A et M le rapport est $- 2,5$.

Bonne réponse : \textbf{Proposition 3}.
\end{enumerate}

\needspace{10\baselineskip}
